%  LaTeX support: latex@mdpi.com 
%  For support, please attach all files needed for compiling as well as the log file, and specify your operating system, LaTeX version, and LaTeX editor.

%=================================================================
\documentclass[water,article,submit,pdftex,moreauthors]{Definitions/mdpi} 

%=================================================================

% MDPI internal commands
\firstpage{1} 
\makeatletter 
\setcounter{page}{\@firstpage} 
\makeatother
\pubvolume{1}
\issuenum{1}
\articlenumber{0}
\pubyear{2024}
\copyrightyear{2024}
\externaleditor{Academic Editor: Firstname Lastname}
\datereceived{7 March 2024} 
\daterevised{11 April 2024} 
\dateaccepted{} 
\datepublished{} 
\hreflink{https://doi.org/} % If needed use \linebreak
%\doinum{}

%=================================================================
% Add packages and commands here. The following packages are loaded in our class file: fontenc, inputenc, calc, indentfirst, fancyhdr, graphicx, epstopdf, lastpage, ifthen, lineno, float, amsmath, setspace, enumitem, mathpazo, booktabs, titlesec, etoolbox, tabto, xcolor, soul, multirow, microtype, tikz, totcount, changepage, attrib, upgreek, cleveref, amsthm, hyphenat, natbib, hyperref, footmisc, url, geometry, newfloat, caption

\graphicspath{{figures/}} % path to folder with figures
\usepackage{subcaption}
\usepackage{listings}
\usepackage{xcolor}
\usepackage{gensymb} % for \textdegree
\usepackage{tabularx}
\newcolumntype{Y}{>{\centering\arraybackslash}X}
\usepackage{amsmath,mathtools,amsthm}
	\setcounter{MaxMatrixCols}{15}
\usepackage{array}
\usepackage{amssymb}
\usepackage{amsfonts}
\usepackage{float}
\usepackage{chngcntr}
\counterwithin*{equation}{section}
\usepackage[super]{nth}
\usepackage{longtable}
\usepackage{tabularx}
\usepackage{multirow}
\usepackage{adjustbox}

\definecolor{deepblue}{rgb}{0,0,0.5}
\definecolor{deepred}{rgb}{0.6,0,0}
\definecolor{deepmagenta}{RGB}{195,12,214}
\definecolor{deepgreen}{RGB}{29,122,30}
\definecolor{mintgreen}{RGB}{192,249,192}
\definecolor{codepurple}{RGB}{204, 0, 153}
\definecolor{LightCyan}{rgb}{0.88,1,1}
\definecolor{GainsboroPurple}{RGB}{247,176,247}
\definecolor{LightSlateGrey}{RGB}{124,118,118}
%    
\lstdefinestyle{mystyle}{
    commentstyle=\color{LightSlateGrey},
    keywordstyle=\color{deepgreen}\bfseries,
    emphstyle=\color{deepred},
    numberstyle=\tiny\color{deepblue},
    stringstyle=\color{codepurple},
    basicstyle=\ttfamily\scriptsize,
    breakatwhitespace=false,         
    breaklines=true,                 
    captionpos=b,                    
    keepspaces=true,                 
    numbers=left,                    
    numbersep=5pt,                  
    showspaces=false,                
    showstringspaces=false,
    showtabs=false,                  
    tabsize=2
}
\lstset{style=mystyle}

%=================================================================
% Full title of the paper (Capitalized)
\Title{Deep Learning Methods of Satellite Image Processing for Monitoring Flood Dynamics in the Ganges Delta, Bangladesh}

% MDPI internal command: Title for citation in the left column
\TitleCitation{Deep Learning Methods of Satellite Image Processing for Monitoring Flood Dynamics in the Ganges Delta, Bangladesh}

% Author Orchid ID: enter ID or remove command
\newcommand{\orcidauthorA}{0000-0002-5759-1089} % Polina Add \orcidA{} behind the author's name

% Authors, for the paper (add full first names)
\Author{Polina Lemenkova $^{1}$*\orcidA{}}

%\longauthorlist{yes}

% MDPI internal command: Authors, for metadata in PDF
\AuthorNames{Polina Lemenkova}

% MDPI internal command: Authors, for citation in the left column
\AuthorCitation{Lemenkova, P.}
% If this is a Chicago style journal: Lastname, Firstname, Firstname Lastname, and Firstname Lastname.

% Affiliations / Addresses (Add [1] after \address if there is only one affiliation.)
\address{%
$^{1}$ \quad Department of Geoinformatics, Faculty of Digital and Analytical Sciences, Universität Salzburg, Schillerstraße 30, A-5020 Salzburg, Austria.
}

% Contact information of the corresponding author
\corres{Correspondence: polina.lemenkova@plus.ac.at; Tel.: +43-677-6173-2772 (P.L.)}

% Current address and/or shared authorship
%\firstnote{Current address: Affiliation 3.} 

% Abstract (Do not insert blank lines, i.e. \\) 
\abstract{Mapping spatial data is essential for monitoring flooded areas, prognosis of hazards and prevention of flood risks. Ganges River Delta, Bangladesh, is the world's largest river delta prone to floods which impact social-natural systems through losses of lives, damage to infrastructure and landscapes. Millions of people living in this region are vulnerable to repetitive floods due to exposure, high susceptibility and low resilience. Cumulative effects from monsoon climate, repetitive rainfalls and tropical cyclones, and hydrogeologic setting in Ganges River Delta increase probability of floods. 
While engineering methods of flood mitigation include practical solutions (technical construction of dams, bridges, hydraulic drains), regulating traffic, land planning support systems, the geoinformation methods relies on modelling remote sensing (RS) data for evaluating the dynamics of flood hazards. Geo-information is indispensable for mapping catchment of flooded areas and visualization of affected regions in real time flood monitoring, implementing and developing emergency plans and vulnerability assessment through warning systems supported by RS data. With this regards, this study uses RS data to monitor southern segment of Ganges River Delta. The multispectral satellite Landsat 8-9 OLI/TIRS satellite images were evaluated in flood (March) and post-flood (November) periods for analysis of flood extent and landscape changes. Deep Learning (DL) algorithms of GRASS GIS and modules of qualitative and quantitative analysis were used as advanced methods of satellite image processing. The results presented a series of maps based on the classified images for monitoring floods in Ganges River Delta.}

% Keywords
\keyword{machine learning; deep learning; flood; monsoon; cartography; geoinformatics; remote sensing; satellite image; geospatial analysis; GRASS GIS} 

\PACS{91.10.Da; 91.10.Jf; 91.10.Sp; 91.10.Xa; 96.25.Vt; 91.10.Fc; 95.40.+s; 95.75.Qr; 95.75.Rs}
\MSC{86A30; 86-08; 86A99; 86A04}
\JEL{Y91; Q20; Q24; Q23; Q3; Q01; R11; O44; O13; Q5; Q51; Q55; N57; C6; C61}

\begin{document}

%----------- section ----------->
\section{Introduction}

\subsection{Background}
\hl{The increasing interest in hydrologic scientific community in flood hazards has dramatically increased the demand for effective methods of geospatial data processing aimed at and detecting floods. }In hydrological studies and\hl{ }management \hl{of coastal regions}, detecting areas prone to floods is essential issue. \hl{Monitoring Earth observation data for }cartographic visualization of flooded landscapes\hl{ supports} flood hazard assessment. \hl{This} includes t\hl{processing Remote Sensing (RS) data for} analysis of \hl{flood extent, consequences of disastrous events, }the degree of flood impact, evaluation of the affected area and \hl{the }estimation of changes in the post-flood periods \cite{su142114145,w15152802,su142316270}. Such analysis enables to highlight\hl{ }spatial extent of\hl{ }flood and \hl{evaluate }spatio-temporal dynamics. \hl{Processing RS data by Geographic Information System (GIS)} has been acknowledged as one of the most powerful approaches in hydrological hazard monitoring. Indeed, \hl{analysis of }satellite images\hl{ }enables to reveal information \hl{on} the consequences of flood \hl{that lead to the } environmental changes such as\hl{ }destroyed areas, \hl{inundated areas}, affected land patches, declined vegetation coverage and disrupted ecosystems \cite{ijgi12110465,rs15102561}. 

Specifically for coastal Asian countries affected by \hl{global climate and environmental changes as well as regional effects from }the monsoon seasonality, fieldwork to capture data on flooded areas is difficult\hl{. Moreover, direct observations in the areas affected by floods can be impossible }both financially and  \hl{practically. Depending} on the safety of the inundated area and severity of hydrological hazards, \hl{onsite surveys }might \hl{be} impossible. Therefore, \hl{using }RS data\hl{ }for mapping flooded landscapes and \hl{evaluating }spatial dynamics of water extent \hl{presents essential }tools for monitoring \hl{floods} \cite{YASEEN2024101665,GHALEHTEIMOURI2023}. 

\subsection{Related works}

Evaluating the degree and rates of floods \hl{is useful for environmental risk assessment. }For hazard prediction and mitigation of hydrological risks\hl{, processing RS data presents the effective approach which can be used to reveal the affected areas. More specifically, }satellite images\hl{ present a} valuable sources of information which can be \hl{processed for detecting areas in pre- and post-flood periods. For example, }they can\hl{ }be used to reveal climate-environmental patterns of floods \hl{and} to evaluate spatio-temporal dynamics in pre- and post-flood periods. To mention a few \hl{examples of such applications}, satellite images have long been used for\hl{ }numerical evaluation of risks \citep{THAKUR2024130683}, hydrological and irrigation modelling \cite{SUN2024415}, susceptibility mapping of the coastal areas \cite{RAZAVITERMEH2023100595}, detecting affected areas and estimating social-economic costs of floods \cite{ZHENG2023101410,DUAN2024131010}, hydrogeological modelling \cite{BELL2021107451}, land use mapping and evaluating climate and environmental impacts on vegetation \cite{JARAMILLO2024526} or computing inundated areas \cite{ZUO2024}.

Processing RS data to retreave information on floods from the satellite images is possible \hl{by} diverse \hl{software and hydrological applications in geo-informatics}. \hl{Traditionally, RS data have been processed by GIS methods.} Diverse tools included in GIS allow mapping the degree and complexity of flood hazards and environmental consequences using classification of the satellite images. For instance, for coastal regions in deltaic areas prone to floods, \hl{the} time series of RS data \hl{are useful} for visualising dynamics of the ecological patterns related to flood events and monsoon cycles \citep{RAZAVITERMEH2023100595,SINGHA2023316,MATHEW2021100861}. Moreover, the analysis of satellite images covering countries located close to shoreline can help to better understand the effects from climate changes on coastal landscapes\hl{.}  

\hl{However, state-of-the-art software severely impedes the timely evaluation and analysis of RS datasets, due to the manual user-adjusted processing techniques. In view of this, }novel \hl{methods} of machine learning (ML) \hl{present better perspectives to image processing and RS data analysis.} ML \hl{emerged recently} as a fundamental technique of image processing and classification\hl{. Its current applications include} a wide variety of \hl{topics such as }environmental studies \cite{ROCHA2022107415}, Earth science \cite{LI2023103345,Lemenkova202306,XIONG2021145256}, geological and hazard risk assessment \cite{AOURAGH2023100939}, hydrological monitoring \cite{RAFIK2023101569} or vegetation analysis \cite{WU2023110723,Lemenkova202310,DOSSANTOS2022106753}\hl{, to mention a few of them}. Technical advantages of ML have been noted earlier \cite{9883389,coasts4010008} \hl{and} include high sensitivity \hl{in data analysis}, flexibility of image \hl{processing} \hl{and accuracy in detecting} variations in \hl{land cover types} to evaluate environmental \hl{dynamics}. \hl{Moreover, diverse parameters of ML techniques} can be finely adjusted using\hl{ different }algorithms \hl{such as} Random Forest (RF), Support Vector Machines (SVM), \hl{and many more}.

\subsection{Gap and motivation}

\hl{In order to perform environmental monitoring of flooded areas, policy makers and scientists need effective techniques and approaches. Therefore, }flood mitigation measures include diverse\hl{ }methods \hl{and tools}. Engineering solutions include, for instance, \hl{the }construction of infrastructure, dams, and bridges, maintaining the hydraulic performance of drains. Regulating traffic based on geo-information regarding \hl{the }location of buildings and land planning support systems enables to support affected communities. The non-engineering methods mostly include modelling, prediction and prognosis\hl{ }for implementing and developing \hl{the }emergency plans. Those enable to assess the vulnerability of area to flood hazards through warning systems supported by \hl{the }up-to-date geo-information. With this regards, using\hl{ }RS data processed by GIS is indispensable for evaluating catchment \hl{of the inundated} areas and visualization of the affected regions in real time flood monitoring.

Currently, the cartographic monitoring of flood in Ganges River Delta \hl{is mostly based on} GIS-based mapping \cite{Sanyal2023,ISLAM2023167,RANA2023}. Besides the conventional approaches, recent studies in the same area study area use alternative \hl{data} such as Sentinel-1 SAR images ad Google Earth Engine \cite{ChakmaAkter,SINGHA2020278}. \hl{Such case studies} raised the question of the operative and accurate methods for data processing aimed at flood prediction and management. Accurate classification of the satellite images for environmental mapping requires advanced methods of scripting languages which aim at robust detecting of the patterns \cite{YANG2020100413,Lemenkova202322,BEVERIDGE2020104843,SADIQ2023111233,Lemenkova202227}. Existing state-of-the-art methods of image processing and classification mostly use the traditional GIS for data handling which may lead to misclassification and inaccurate labelling of pixels. \hl{Nevertheless, the software that process RS data for monitoring floods must include intelligent algorithms and techniques to deal with complex spatial patterns of the flooded coastal areas, i.e., advanced computational methods that exploit spatial variability and heterogeneity of the affected landscapes, rather than visualise particular extents of the images.} \hl{With this regards, ML and programming} techniques \hl{present advanced solutions} to support real-time flood prognosis systems through the algorithms of artificial intelligence\hl{ (AI)} \cite{ZHAO2021126777}. \hl{Such methods can be used for operative monitoring of floods and }identifying landscape patterns within coastal and tidal areas, estuary and delta \hl{which }is challenging \hl{task }due to the obscurity and ambiguity of land cover types. 

\hl{Complimentary to} the available traditional tools of satellite image\hl{ }classification, novel methods of Deep Learning (DL) can now be used as \hl{an} advanced approach \hl{for} RS \hl{data processing}. For instance, such algorithms as Multilayer Perceptron (MLPClassifier) are designed to support image partition, classification and analysis through algorithms of decision trees and computer vision \cite{6493388,6422192,7507955}. The application of such methods in hydrological studies and coastal risk assessment creates principally new perspectives in cartography by combining RS \hl{data }with programming approaches for mapping flooded areas. It is therefore particularly important to apply such advanced technical tools to regions with high heterogeneity and complexity of landscape patterns, such as coastal zones, where algorithms of DL can generate decisions on image classification and detect affected areas \hl{automatically}. 

When applied to a series of \hl{the }satellite images, ML methods can effectively discover the properties of landscapes through comparative analysis \cite{9324433,Lemenkova202313}. Moreover, ML enables to quantify the patches of landscapes on raster images using algorithms of computer vision  \cite{10281789,jmse11040871,8899013,Lemenkova202021,9333522} or determine landscape dynamics for environmental monitoring \cite{9392684,Lemenkova202324}. Finally, other advantages of the ML approach is that it is a resource- and time-effective method based on advanced programming methods which largely involves scripting techniques. Such approach enables to automate data processing and cartographic analysis as mentioned earlier \cite{10282693,Lemenkova202301,10441504}. Hence, using ML enables to indicate flood risk and development of crisis situation in flood-prone areas such as delta of Ganges, Bangladesh. Moreover, such methods are especially effective for detecting correlation between the environmental processes \hl{in coastal areas }and \hl{effects of }climate change over time which is especially actual for tropical regions that are subject to monsoon \hl{activity}. \hl{Nevertheless, among the most serious deficiencies scientists face with current GIS is the lack of appropriate ML-based methods of RS data processing. With this regards, Geographic Resources Analysis Support System (GRASS) GIS presents an advantageous software that includes methods of ML for manipulating the satellite images using advanced algorithms such as Deep Learning (DL) and Artificial Neural Network (ANN).}

\subsection{Goals and objectives}

This paper presents the use of the advanced DL tools of GRASS GIS for \hl{satellite }image classification. \hl{Moreover, it addresses a particular problem of monitoring flooded areas around Ganges River Delta, Bangladesh.} Several Landsat satellite images are used, processed, classified, compared and analysed to evaluate \hl{the }pre-and post flooded periods in coastal Bangladesh. The potential of the ML modules of GRASS GIS for mapping consequences of \hl{the }hydrological hazards is exemplified with the case of \hl{the }affected landscapes of coastal Bangladesh and associated climate dynamics related to the monsoonal climate of the Indian Ocean. The focus of the manuscript is on a holistic approach to accurate satellite image processing with the help of the open source software GRASS GIS. The methodology applied to the subject of the study can be applied to the other relevant countries located in coastal regions with flood-prone areas.

\section{Study Area}

\subsection{Current landscape}

This study presents a case of the Ganges River Delta, \hl{the area regularly affected by floods and }located in southern Bangladesh\hl{, }Figure \ref{fig01}. 

\begin{figure}[H]
    \begin{center}
	\hspace{-35pt}\includegraphics[width=11.0cm]{Fig_01.jpg}
    \end{center}
    \caption{Topographic map of Bangladesh with depicted study area showing the placement of the Landsat data over the Ganges River Delta. Digital elevation data: SRTM/GEBCO, 15 arc sec resolution grid \cite{GEBCO,Tozer}. Software: GMT version 6.4.0 \cite{Wessel}. Map source: author.}
    \label{fig01}
\end{figure}

The Ganges River Delta is the\hl{ }largest river delta \hl{in the world. It includes }hundreds of millions people living around in a total catchment area. Located in the region with monsoon climate and seasonally changing environmental setting, it is subject to \hl{regularly repeating} floods, occasional tropical cyclones, tidal waves and related inundations \cite{Bhardwaj2022}. Complex hydrology of the Ganges River Delta increases the consequences of floods through dense network of streams, rivers and tributaries where water remains for a longer period \cite{Islam}. Alluvial clayey sediments dominating in the riverbed contribute to the stagnation of water during flood \hl{since it} accumulates in numerous interconnected channels, inner lakes, wetlands and flood plains of the Ganges River Delta \cite{Datta,KHAN201927,Lemenkova202029}. The monsoon climate regulates the repeatability of rainfalls and floods in Ganges River Delta which normally \hl{increase in the period} between\hl{ }March and September. 

\subsection{Problem formulation}
The repetitive floods in Ganges River Delta negatively impact social-natural systems\hl{. They cause }losses of lives and damage to infrastructure \cite{JERIN2023}, fragmentation and disruption in natural landscapes \cite{MUKHERJEE2021106961} and \hl{increase salinity due to intrusion of oceanic water }during storms. \hl{Moreover, }floods affect crop agriculture \cite{ABEDIN2023315,ALMASUD2020104443}, and negatively influence the sustainability of coastal ecosystems \cite{SARKER2024169718,MAINUDDIN2021105740,HASAN2023101028}. Millions of people living in southern Bangladesh and neighbouring regions of India are vulnerable to \hl{such hazards }due to exposure, high susceptibility and low resilience \hl{to repetitive floods}\cite{SINGHA2020106825,RUMPA2023104047}. 

The examples of social effects from floods include the \hl{restricted or limited access} to clean drinking water\hl{, damage }electricity \hl{network}, broken transport and communication systems\hl{. Furthermore, flood hazards disrupt} regular education and normal working environment \hl{as well as} health care support for population \cite{JERIN2023103851,NAHIN2023e14520}. The long-term effects from \hl{floods and their} consequences which may last for months  seriously affect the social-economic system of Bangladesh. The coastal area of Ganges River Delta supports approximately 56.7 M of population which encompasses the study area \cite{DAS2021101983} where many settlements are located\hl{. This includes }both small villages and large cities such as the capital Dhaka, Khulna (the third-largest city in the country), Comilla and Narayani. Supporting such flood-prone areas requires geospatial monitoring of flood risks in order to visualise the affected areas on the maps of floods \cite{Alietal,AZAD2023100498}.

The uncertainty of flood prediction in Ganges River Delta has an impact on vulnerability to risk and information among the local population. Hence, preventive prediction rises a question of using reliable data and advances methods for operative monitoring of floods and assessment of environmental sustainability \hl{using flood maps} \cite{NICHOLLS2016370}. \hl{Various} techniques have been developed to produce flood maps. Among them, satellite-based flood maps tend to provide real-time, contiguous representations with varying spatial resolution over time and higher accuracy. The main approach that a country can take to assess flood risk is to map areas that are susceptible to flooding based on their \hl{geographic setting} and\hl{ }comprehensive flood risk assessment. \hl{Such analysis can be performed when satellite images are integrated with geospatial information, e.g., frequency of floods, topographic relief, population density, location of cities and climate setting. }With this regards, flood maps should be prepared by the relevant authorities, and remote flood monitoring should be used to reduce the risk during weather events. 

Notably, no prior study has explored the deep learning (DL) methods available in \hl{GRASS GIS} \cite{GRASS_GIS_software}\hl{ }for mapping floods in Ganges River Delta, Bangladesh. As a response to this gap, this study presents \hl{a comprehensive evaluation of multi-spectral imagery as }a report with such approach\hl{. Specifically, the methodology of this research }is based on using the DL methods of satellite image processing by GRASS GIS scripting software. Deep learning is highly scalable for image processing purposes since it can analyse large datasets such as times series of the RS data and conducts numerous computations in a cost- and time-effective way using scripting software such as GRASS GIS. Hence, the application of DL for environmental monitoring increases productivity of RS data processing by speeding up image processing using its different algorithms, Figure \ref{fig02}. 

\begin{figure}[H]
    \begin{center}
%    \setlength{\unitlength}{0.5cm}
	\hspace{-35pt}\includegraphics[width=8.0cm]{Fig_02.jpg}
    \end{center}
    \caption{Conceptual structure of the ML and DL algorithms in geospatial data analysis and image processing by GRASS GIS. Graphical software: Inkscape version 1.2 \cite{Inkscape}. Scheme source: author.}
    \label{fig02}
\end{figure}

Moreover, DL\hl{ }also increases\hl{ the }flexibility of\hl{ }image processing by allowing trained models\hl{ }applied to various environmental \hl{and hydrological }problems. The implementation of \hl{such} approach is realised through the MLPClassifier algorithm from the embedded Scikit-Learn library of Python \cite{pedregosa2011scikit,van1995python}. 

%----------- section ----------->
\section{Materials and Methods}

%----------- subsection ----------->
\subsection{Deep Learning}

As a branch of Machine learning (ML), Deep learning (DL) is increasingly used in \hl{the }Earth \hl{sciences as a tool} for automatic processing of geospatial data. \hl{The }DL can be applied by a variety of tools, using programming libraries, embedded plugins and modules in Geographic Information Systems (GIS). \hl{However, the} problem remains with the technical approach of such methods \hl{specifically for RS data and satellite image processing}. Amongst the geoinformation software \hl{that operate with RS data}, GRASS GIS \cite{GRASS_GIS_software,NETELER2012124,MitasovaNeteler} proposes the effective solutions to satellite image processing with the \hl{valuable }option \hl{of utilising} diverse choices of models for image classification.

\hl{Among} the existing ML methods, the most widely used approach is the DL with many examples of applications for \hl{satellite} image processing in environmental studies \cite{10441152,AHMED2024167234,HASSAN2023100399}. \hl{For example, the }DL presents the advanced technique for evaluating the environmental dynamics during climate change assessed from the satellite image analysis \cite{10441576}. There are many DL algorithms which have various parameters and functionality for image classification. Overall, \hl{such algorithms} train the computer model using knowledge derived from the input layers and processed in hidden layers which results in accurate modelling of data. \hl{Moreover, the DL considers existing conditions in geospatial data which are commonly expressed in terms of spatial relations and topology of objects.}

Existing algorithms of the ML modules of GRASS GIS use Neural Networks (NN) \cite{202017,HECHTNIELSEN199265,SHEKHAR199213} to learn semantic representations of features as indicated by pixels on the images derived from Earth Observation (EO) data. This enables to evaluate landscapes dynamics  and \hl{assess }environmental properties \cite{BISWAS2023e21245}\hl{, apply advanced methods of} segmentation \cite{Lemenkova202320}, feature extraction, and classification \cite{RAISA2024101128} \hl{for data analysis}. Further development of \hl{the }DL methods aims at achieving high precision of\hl{ }image processing as \hl{information} obtained from spectral reflectances of the land cover types detected on the images, as well as the use of complex programming algorithms, e.g., NN \cite{RUDRA2023e16459}. 

%----------- subsection ----------->

\subsection{Workflow}
The methodological workflow consists of several steps aimed at geospatial analysis and image processing\hl{, }Figure \ref{fig03}. 

\begin{figure}[H]
    \begin{center}
	\includegraphics[width=10.0cm]{Fig_03.jpg}
    \end{center}
    \caption{Workflow scheme \hl{illustrating} methodology used in this study for image processing \hl{by} GRASS GIS. Software: R, library DiagrammeR \cite{DiagrammeR}. Diagram source: author.}
    \label{fig03}
\end{figure}

\hl{This} study \hl{mainly uses} the DL methods available in the GRASS GIS modules for satellite image processing to detect flooded areas in Ganges River Delta, Bangladesh. \hl{The methodology} also included \hl{processing }diverse data obtained from different sources. While the purpose of \hl{the }EO data \hl{was} to provide information on objects and features \hl{of land surface }visible from space, the choice of methods serves varied \hl{according to }purposes that \hl{include} specific DL algorithms and techniques of image processing \cite{4736879,5070195}. 

%----------- subsection ----------->
\subsection{Theoretical fundamentals of MLP}
The Multilayer perceptron (MLP) is a class of DL and a conceptual theoretical background to data analysis performed in this study. MLP presents a feedforward Artificial Neural Network (ANN), a branch of the DL methods \cite{su151813786}. It consists of fully connected neurones with a nonlinear activation function organised in at least three layers, Figure \ref{fig04}. 

\begin{figure}[H]
    \begin{center}
	\includegraphics[width=14.0cm]{Fig_04.jpg}
    \end{center}
    \caption{General methodological scheme for \hl{the }DL \hl{method of} satellite image classification and EO data processing. Software: R, library DiagrammeR \cite{DiagrammeR}. Diagram source: author.}
    \label{fig04}
\end{figure}

These layers are used to distinguish \hl{the }data that are not linearly separable and, for time series of satellite images,\hl{ }variables that are changing over time. This results in different values of pixels recognised by the model as Digital Numbers (DN). Such characteristics is caused by\hl{ }spatio-temporal variability of spectral reflectance of pixels that constitute the image\hl{. Since }physical properties of landscape and vegetation patches \hl{vary significantly, such mosaic of landscape patches is reflected on the raster matrix and can be identified on the satellite images, accordingly} \cite{rs16010184,HUANG20181915,YIN2021102514}. 

\subsection{Data}
In this study, we used eight Landsat 8-9 Operational Land Imager (OLI) and Thermal Infrared Sensor (TIRS) multispectral satellite images for environmental mapping of Ganges Delta, Bangladesh. \hl{The} images are selected for flood and post-flood periods (March and November) \hl{covering periods of }2021, 2022 and 2023. The images \hl{on} March and November of\hl{ }2014 were used as a seed of training pixels for Deep Learning classification. The generated dataset is collected from the \href{https://earthexplorer.usgs.gov/}{EarthExplorer} \cite{USGSPublication} with image samples openly available for download, processing and data reuse. The original Landsat images included multispectral, panchromatic and SWIR bands, Figure \ref{fig05}. 

The images are in GeoTIFF format and contained the seven multispectral bands for each image. Hence, the image recordings received from the USGS source include RS datasets\hl{. The metadata were} explored \hl{and the} records \hl{are }available in the Appendix in Table \ref{tabApp}\hl{. The most }essential \hl{technical }characteristics \hl{of the images }are summarised in Table \ref{tab01}. \hl{Besides quality, coverage and availability, another }advantage of the Landsat archives \hl{consists in their economic values: these data are reusable for further distribution and future research}. Thus, in \hl{similar} studies, these \hl{scenes} can be \hl{utilised} as input data for similar tasks on \hl{image analysis}, such as classification\hl{, }detection of land cover types, and image segmentation\hl{. The }presented GRASS GIS workflow model of other \hl{technical methods described in this study can be utilised for image processing}. 

The datasets are available in EarthExplorer \cite{USGSPublication} and \hl{originate }from \hl{the }NOAA Landsat collections \cite{Landsat}. The images were pre-projected as original data and had the following technical characteristics: Datum and Ellipsoid World Geodetic System 84 (WGS84), stored in Universal Transverse Mercator (UTM) Zone 46 for southern Bangladesh, and Worldwide Reference System (WRS) Path/Row 137/44. The data were collected at Nadir in day time. The Landsat Collection Category \hl{is} T1, Number 2 for all the scenes. \hl{The} Landsat Station Identifier is LGN \hl{with }sensor Identifier \hl{as} Landsat OLI TIRS for all the images\hl{. The }Data Type L2 is OLI TIRS L2SP. The remaining essential characteristics of the images are summarised in Table \ref{tab01}.

\begin{table}[H]
%\scriptsize
\caption{Attribute table of the main characteristics of the Landsat 8-9 OLI/TIRS images} \label{tab01}
%\begin{adjustwidth}{-\extralength}{0cm}
 %   \begin{tabularx}{\fulllength}{llll}
 %    \begin{tabularx}{\textwidth}{cccc}
\begin{tabularx}{\textwidth}{@{}lYYYYY@{}}
    \toprule
       	\textbf{Data Set Attribute} & \textbf{Attribute Value} & \textbf{Attribute Value} & \textbf{Attribute Value} \\ \hline
        \textbf{Date Acquired} & \textbf{2023/03/07} & \textbf{2022/03/20} & \textbf{2021/03/17} \\ \hline
        \textbf{Land Cloud Cover} & 3.12 & 0.01 & 0.03 \\ \hline
        \textbf{Scene Cloud Cover L1} & 2.97 & 0.01 & 0.03 \\ \hline
        \textbf{Sun Elevation L0RA} & 51.70012312 & 56.10505202 & 55.19368850 \\ \hline
        \textbf{Sun Azimuth L0RA} & 134.86314148 & 129.90704446 & 131.01649112 \\ \hline
        \textbf{Satellite} & 8 & 8 & 8 \\
        \midrule
        \textbf{Data Set Attribute} & \textbf{Attribute Value} & \textbf{Attribute Value} & \textbf{Attribute Value} \\ \hline
        \textbf{Date Acquired} & \textbf{2023/11/26} & \textbf{2022/11/23} & \textbf{2021/11/28} \\ \hline
        \textbf{Land Cloud Cover} & 0.02 & 0.20 & 0.40 \\ \hline
        \textbf{Scene Cloud Cover L1} & 0.02 & 0.19 & 0.38 \\ \hline
        \textbf{Sun Elevation L0RA} & 41.83496249 & 42.43660966 & 41.36291892 \\ \hline
        \textbf{Sun Azimuth L0RA} & 154.44654325 & 154.42300606 & 154.52745495 \\ \hline
        \textbf{Satellite} & 9 & 9 & 8 \\
        \bottomrule
    \end{tabularx}
%    \end{adjustwidth}
\end{table}

The selection of these data is explained by the reputation and reliability of the Landsat \hl{satellite }imagery widely used for environmental studies. Besides, the Landsat images are \hl{openly }available free of charge, have high \hl{frequency of orbit (8-day repeat coverage of Landsat scene area)} and high quality. Hence, \hl{these }images present a reliable source of information for detecting the extent of floods and visualising inundated areas using time series analysis. Their applications is reported in many existing studies with a focus on environmental monitoring \cite{Sheonty,Islam2016,Lemenkova202229,Ferdous,Lemenkova202321}. The data were collected from the \hl{freely} available  \hl{repository of }the United States Geological Survey (USGS)\hl{ -- the }EarthExplorer \cite{USGSPublication}. The original data are visualised in Figure \ref{fig05}. 

\begin{figure}[H]
  \centering
  \begin{tabular}[b]{c}
    \includegraphics[width=2.9cm]{Fig_05a.jpg} \\
    \small (a): March 2014
  \end{tabular}
  \begin{tabular}[b]{c}
    \includegraphics[width=2.9cm]{Fig_05b.jpg} \\
    \small (b): November 2014
  \end{tabular}
    \begin{tabular}[b]{c}
    \includegraphics[width=2.9cm]{Fig_05c.jpg} \\
    \small (c): March 2021
  \end{tabular}
  \begin{tabular}[b]{c}
    \includegraphics[width=2.9cm]{Fig_05d.jpg} \\
    \small (d): November 2021
  \end{tabular}
   \begin{tabular}[b]{c}
    \includegraphics[width=2.9cm]{Fig_05e.jpg} \\
    \small (e): March 2022
  \end{tabular}
  \begin{tabular}[b]{c}
    \includegraphics[width=2.9cm]{Fig_05f.jpg} \\
    \small (f): November 2022
  \end{tabular}
    \begin{tabular}[b]{c}
    \includegraphics[width=2.9cm]{Fig_05g.jpg} \\
    \small (g): March 2023
  \end{tabular}
  \begin{tabular}[b]{c}
    \includegraphics[width=2.9cm]{Fig_05h.jpg} \\
    \small (h:) November 2023
  \end{tabular}
  \caption{Data sources: Landsat 8-9 OLI/TIRS images on Ganges River Delta, Bangladesh, collected during the flood period (March) and post-flood period (November) on years 2014, 2021, 2022 and 2023, used for analysis of flood dynamics. Data source: USGS \cite{Landsat}. Compilation source: author.}
  \label{fig05}
\end{figure}

\subsection{Software}
The workflow used in this research employs a chain of modules that \hl{utilise} scripting GRASS GIS software \cite{GRASS_GIS_software} for image processing\hl{. The technical details of such approach are }described in earlier works \cite{Lemenkova202323,technologies11020046}. For the installation and use of Machine Learning (ML) modules of GRASS GIS, we recommend using a Python version 3.8.5 or higher. Specifically, the installation of our framework can be performed using the available packages in the shared data. The additional Python libraries required by GRASS GIS are included in the modules through \hl{the} embedded links and include the following \hl{ones}: Matplotlib, NumPy, Pandas, Scikit-Learn, SciPy. Additionally, the XCode \hl{needs} to be installed for writing the scripts\hl{.}  These libraries can be installed using 'pip install' with more details in the online resources of the GRASS GIS.  \hl{Finally, }using GitHub is recommended \hl{for backup of scripts and revision}.The methodology of \hl{the} time series analysis includes the unsupervised and supervised classification of \hl{the} multispectral satellite images \hl{aimed at} detecting changes in landscapes around the Ganges River Delta in \hl{the} pre- and post-flood periods. The map in Figure \ref{fig01} is prepared using Generic Mapping Tools (GMT) version 6.4.0 \cite{Wessel,WesselSmith} using raster grid of General Bathymetric Chart of the Oceans (GEBCO). 

%----------- subsection ----------->
\subsection{Implementation}
We used a workflow in GRASS GIS platform as a container of diverse modules to run the scripts on MacOS using the approach presented in Figure \ref{fig05}. The GRASS GIS uses algorithms of ML/DL methods of image processing derived from the Python's Scikit-Learn library \cite{pedregosa2011scikit} via the integrated libraries in the machine. For GRASS GIS platform, this is performed using the containers of modules via the sequence of commands presented below in scripts in Listings \ref{lst01}, \ref{lst02} and \ref{lst03}. The processes include the import and pre-processing of the multispectral bands of the given image, image analyse and classification\hl{. The classification is based on fundamental properties of the RS data of }difference in spectral reflectance \hl{of pixels that are }associated with land patches \hl{on the land surfaces visible from space. Here, }each\hl{ }patch of the Ganges River Delta \hl{contributes to the mosaic of }landscapes evaluated for March (flood) and November (post-flood) periods of 2021, 2022 and 2023.

The following is an outline of the essential details regarding the approaches and techniques of the algorithm that compose the GRASS GIS framework \hl{used} for image analysis. For details, full codes are presented below and explained with \hl{the} essential details. In all the GRASS GIS modules and functions and post-processing of the satellite images using DL approach, it is necessary to apply the module which \hl{implements} the required functionality for data processing for each image\hl{. In this sense, the parameters of scripts are adjusted }in terms of the identifiers of functions and \hl{specifications} of image processing. For this purpose, we adopted and modified the approaches of ML modules\hl{ }of GRASS GIS which rely on\hl{ }the Python's Scikit-Learn library for DL techniques.

%----------- subsection ----------->
\subsection{Image processing}
The images were first imported into the GRASS GIS system using 'r.import' module\hl{ and pre-processed. Afterwards, the }colour composites \hl{were} created using 'r.composite'. The images were then displayed using the combination of modules 'd.mon' and 'd.rast'. The files \hl{were} saved to the bitmap format using the 'd.out.file' module. The binary results of the DL mapping and programming scripts used for DL image analysis are available in the GitHub repository, \hl{using} the following web link: \\
\href{https://github.com/paulinelemenkova/Floods_Ganges_GRASS_GIS_DL_Image_Analysis}{https://github.com/paulinelemenkova/Floods\_Ganges\_GRASS\_GIS\_DL\_Image\_Analysis}.

The results of the \hl{generated }colour composites are demonstrated in Figure \ref{fig06}. 

\begin{figure}[H]
	\begin{subfigure}[b]{.45\textwidth}
		\centering
		\includegraphics[width=0.98\textwidth]{Fig_06a.jpg}
		\caption{Bands 3-4-7}
	\end{subfigure}%
	\begin{subfigure}[b]{.45\textwidth}
		\centering
		\includegraphics[width=0.98\textwidth]{Fig_06b.jpg}
		\caption{Bands 3-4-5}
	\end{subfigure}
		\hfill
		\vspace{1em}
	\begin{subfigure}[b]{.45\textwidth}
		\centering
		\includegraphics[width=0.98\textwidth]{Fig_06c.jpg}
		\caption{Bands 2-4-7}
	\end{subfigure}
	\begin{subfigure}[b]{.45\textwidth}
		\centering
		\includegraphics[width=0.98\textwidth]{Fig_06d.jpg}
		\caption{Bands 3-7-4}
	\end{subfigure}
\caption{Examples of different false colour composites generated using the multispectral bands of the Landsat 8-9 OLI/TIRS scenes by GRASS GIS 'r.composite' module. Here the band numbers correspond to the following channels of the Landsat 8-9 OLI/TIRS images: Band 2 -- Blue; Band 3 -- Green; Band 4 -- Red; Band 5 -- Near Infrared (NIR); Band 7 -- Shortwave Infrared-2 (SWIR-2).}\label{fig06}
\end{figure}

The implementation of this step is performed using the GRASS GIS syntax (example is presented for November 2023), presented in the codes of Listing \ref{lst01}.

\begin{lstlisting}[language=bash,caption=GRASS GIS code for data import and creating colour composites algorithm,style=mystyle,label={lst01}]
#importing the image subset with 7 Landsat bands and display the raster map
r.import input=/Users/polinalemenkova/grassdata/Bangladesh/LC09_L2SP_137044_20231126_20231128_02_T1_SR_B1.TIF output=L8_2023_N_01 extent=region resolution=region --overwrite
r.import input=/Users/polinalemenkova/grassdata/Bangladesh/LC09_L2SP_137044_20231126_20231128_02_T1_SR_B2.TIF output=L8_2023_N_02 extent=region resolution=region
r.import input=/Users/polinalemenkova/grassdata/Bangladesh/LC09_L2SP_137044_20231126_20231128_02_T1_SR_B3.TIF output=L8_2023_N_03 extent=region resolution=region
r.import input=/Users/polinalemenkova/grassdata/Bangladesh/LC09_L2SP_137044_20231126_20231128_02_T1_SR_B4.TIF output=L8_2023_N_04 extent=region resolution=region
r.import input=/Users/polinalemenkova/grassdata/Bangladesh/LC09_L2SP_137044_20231126_20231128_02_T1_SR_B5.TIF output=L8_2023_N_05 extent=region resolution=region
r.import input=/Users/polinalemenkova/grassdata/Bangladesh/LC09_L2SP_137044_20231126_20231128_02_T1_SR_B6.TIF output=L8_2023_N_06 extent=region resolution=region
r.import input=/Users/polinalemenkova/grassdata/Bangladesh/LC09_L2SP_137044_20231126_20231128_02_T1_SR_B7.TIF output=L8_2023_N_07 extent=region resolution=region
g.list rast
# false color
r.composite blue=L8_2023_N_07 green=L8_2023_N_05 red=L8_2023_N_03 output=L8_2023_N_753
d.mon wx0
d.rast L8_2023_N_753
d.out.file output=Bangladesh_753 format=jpg --overwrite
# false color: NIR band B05 in the red channel, red band B04 in the green channel and green band B03 in the blue channel
r.composite blue=L8_2023_N_03 green=L8_2023_N_04 red=L8_2023_N_05 output=L8_2023_N_345
d.mon wx0
d.rast L8_2023_N_345
d.out.file output=Bangladesh_345 format=jpg --overwrite
# true color
r.composite blue=L8_2023_N_02 green=L8_2023_N_03 red=L8_2023_N_04 output=L8_2023_N_234
d.mon wx0
d.rast L8_2023_N_234
d.out.file output=Bangladesh_J_234 format=jpg --overwrite
\end{lstlisting}

The unsupervised classification is based on clustering \hl{technique} using k-means and Maximum Likelihood embedded in the GRASS GIS. \hl{Practically}, it was performed using the modules 'i.group', 'i.cluster ' and 'i.maxlik'. \hl{The latter creates} the signature file and performs accuracy assessment using \hl{the} 'reject' function\hl{. This function }evaluates \hl{the }rejection probability of the pixels using chi-square test \hl{which is used to examine the dependency of categorical variables, i.e., pixels on the raster scene}. The implementation of this step is performed using the GRASS GIS syntax \hl{with the }example\hl{ }presented \hl{below }for November 2023),\hl{ }in\hl{ }codes of Listing \ref{lst02}.

\begin{lstlisting}[language=bash,caption=GRASS GIS code for unsupervised image classification method using k-means clustering algorithm,style=mystyle,label={lst02}]
g.region raster=L8_2023_N_01 -p
i.group group=L8_2023_N subgroup=res_30m \
  input=L8_2023_N_01,L8_2023_N_02,L8_2023_N_03,L8_2023_N_04,L8_2023_N_05,L8_2023_N_06,L8_2023_N_07 --overwrite
# Clustering: generating signature file and report using k-means clustering algorithm
i.cluster group=L8_2023_N subgroup=res_30m \
  signaturefile=cluster_L8_2023_N \
  classes=10 reportfile=rep_clust_L8_2023_N.txt --overwrite
# Classification by i.maxlik module
i.maxlik group=L8_2023_N subgroup=res_30m \
  signaturefile=cluster_L8_2023_N \
  output=L8_2023_N_cluster_classes reject=L8_2023_N_cluster_reject --overwrite
\end{lstlisting}

The results of this step\hl{ }are used as training data for supervised classification using deep learning approach. \hl{The resulting maps are demonstrated in Figure} \ref{fig07}. Moreover,\hl{ }selected example of accuracy analysis for 2014 which was used as training data set is shown in Figure  \ref{fig08}. This step includes the procedures of clustering that generated signature file and \hl{a report} which was performed using k-means algorithm \hl{by} 'i.cluster' module \hl{of GRASS GIS}. The unsupervised classification was done using the 'i.maxlik' module. The results of this step\hl{ }provided the seed data and training map for the next step of DL modelling\hl{. The maps are presented in Figure} \ref{fig02}.

The results of the accuracy assessment are presented in Figure \ref{fig08} which shows the reject threshold map layer created using GRASS GIS module 'i.maxlik'. The maps are computed for each satellite image and contain the index to one calculated confidence level for each classified pixels in the classified images. In total, the values range from 0 to 100 \% which indicate the predefined confidence intervals. The interpretation of the reject maps includes values of lower (is to be  as 0.1 = keep and those close to 100\% indicate reject. The application for this map layer is as a mask indicating raster pixels in the classified images that have a low probability of flood categorisation (that is, high rejection index) of being assigned to the correct class. 

The computed maps \hl{are} based on the evaluated reject raster map which indicates the results of the evaluated reject thresholds. The reject thresholds were calculated using a chi square test which is a statistical approach that tests \hl{the} goodness of fit of the datasets. This test was used to determine whether the computed pixels and assigned to various land cover classes are acceptable or should be rejected\hl{. The approach of the chi square test is }based on the hypothesis that the resulting data \hl{have} specified distribution of land cover classes at the specified level of significance. Hence, the accuracy assessment was tested on each satellite image and shows the discriminant result at various threshold levels of confidence level of the assigned pixels. The maps shown in Figure \ref{fig08} \hl{are} used to determine at what confidence level each classified pixels is correctly categorised to the diverse land cover classes. 

\begin{figure}[H]
	\begin{subfigure}[b]{.50\textwidth}
		\centering
		\includegraphics[width=0.98\textwidth]{Fig_07a.jpg}
		\caption{March 2014}
	\end{subfigure}%
	\begin{subfigure}[b]{.50\textwidth}
		\centering
		\includegraphics[width=0.98\textwidth]{Fig_07b.jpg}
		\caption{November 2014}
	\end{subfigure}
\caption{Image processing of the Landsat 8 OLI/TIRS scene on March and November 2014 using unsupervised classification by clustering and Maximal Likelihood method: \textbf{(a)}: March 2014; \textbf{(b)}: November 2014. Ten land cover types correspond to the following categogies: 1); Water; 2) Wetlands and mudflats; 3) Mangrove forests; 4) Sandy areas; 5) Forests; 6) Croplands; 7) Grasslands; 8) Urban settlements; 9) Orchards; 10) Aquaculture.}\label{fig07}
\end{figure}

\begin{figure}[H]
	\begin{subfigure}[b]{.50\textwidth}
		\centering
		\includegraphics[width=0.98\textwidth]{Fig_08a.jpg}
		\caption{March 2014}
	\end{subfigure}%
	\begin{subfigure}[b]{.50\textwidth}
		\centering
		\includegraphics[width=0.98\textwidth]{Fig_08b.jpg}
		\caption{November 2014}
	\end{subfigure}
\caption{Accuracy assessment of image classification of the Landsat 8 OLI/TIRS scenes: \textbf{(a)}: March 2014; \textbf{(b)}: November 2014.}\label{fig08}
\end{figure}

\hl{The} Deep learning (DL) analysis is performed using algorithm of Multilayer Perceptron (MLPClassifier) for automated classification of the series of satellite images. First, the spatial region extent was defined using \hl{the} 'g.region' module which sets up the coordinates of the maps within the project using the target map. Afterwards, the training pixels from an older (2014) land cover classification were generated using \hl{the} 'r.random' module of GRASS GIS and used as seeds. Next, the imagery group with all Landsat-8 OLI/TIRS bands was created to include all \hl{the} multispectral bands (\hl{since }we do not need panchromatic bands in this case \hl{they were excluded from data analysis}). Then the training pixels were applied to perform a classification using \hl{the} 'MLPClassifier' of \hl{DL} approach \hl{applied to} recent Landsat \hl{image }(in the presented code below -- the image on November 2023). This method trains a MLPClassifier model using 'r.learn.train. 

After this step, the prediction of the pixel assignments to each land cover class over the Ganges River Delta was performed using \hl{the} 'r.learn.predict' module. The raster categories were automatically applied to the classification output and checked again using the 'r.category' module. \hl{Following that}, the color scheme was assigned from the land class training map to result and \hl{visualised} using the modules 'r.colors' and 'd.rast'. Technically, the DL approach was realised using the 'r.learn.train' module of GRASS GIS and demonstrated in script of Listing \ref{lst03}.

\begin{lstlisting}[language=bash,caption=GRASS GIS code for unsupervised image classification method using k-means clustering algorithm,style=mystyle,label={lst03}]
g.region raster=L8_2023_N_01 -p
r.random input=L8_2014_N_cluster_classes seed=100 npoints=1000 raster=training_pixels
i.group group=L8_2023_N input=L8_2023_N_01,L8_2023_N_02,L8_2023_N_03,L8_2023_N_04,L8_2023_N_05,L8_2023_N_06,L8_2023_N_07 --overwrite
r.learn.train group=L8_2023_N training_map=training_pixels \
    model_name=MLPClassifier n_estimators=500 save_model=mlpc_model.gz --overwrite
r.learn.predict group=L8_2023_N load_model=mlpc_model.gz output=mlpc_classification_Ganges --overwrite
r.category mlpc_classification_Ganges
d.mon wx0
d.rast mlpc_classification_Ganges
r.colors mlpc_classification_Ganges color=roygbiv -e
d.legend raster=mlpc_classification_Ganges title="MLPClassifier: 11/2023" title_fontsize=14 font="Helvetica" fontsize=12 bgcolor=white border_color=white
d.out.file output=MLPC_2023_11 format=jpg --overwrite
\end{lstlisting}

Each image of the produced dataset contained seven multispectral bands processed automatically by GRASS GIS modules 'r.learn.train' and 'r.learn.predict' for each image. The total time processing for one image is ca. 30 minutes, which demonstrates a high time efficiency of the proposed GRASS GIS algorithm of \hl{DL}. The algorithm of MLPClassifier was tested initially on a single image processed as a complete workflow for datasets on March and November 2021, 2022 and 2023 and then analysed for changes in the evaluated period. Changes detected using the repeatability pattern of land cover types over the Ganges River Delta were evaluated for the flood- and post-flood periods. Noise background related to cloudiness is ignored since the images were collected with a high visibility (below 10\% of cloudiness) quality which ensured that land cover types are clearly visible and identified by the ML algorithm. 

%----------- section ----------->
\section{Results}

Satellite images Landsat evaluated in this study cover the regions of southern Bangladesh and periods of March and November for years 2021, 2022 and 2023\hl{, respectively}. \hl{The }analysed images included multispectral bands, and were collected on the cloud-free dates which enabled to detect notable changes in the flooded areas. Land cover changes during flood- and post-flood periods were measured using the advanced 'MLPClassifeir' approach of \hl{DL, as described in previous section. The scripting approach in }developed \hl{in} GRASS GIS which employs the Pythons's algorithms of Scikit-Learn library of \hl{ML}. The variations in the flooding characteristics in the mapped region show the increase of the water level which is associated the pre-monsoon and monsoon rainfall over the river catchments of the Ganges, Brahmaputra and Meghna rivers. The addition factor to the water flow is the intake from the snow melt in February which is reflected in early March levels. 

The hydrological peaks of the Brahmaputra River in late August or September showing the maximum flow of a stream in response to a rainstorm event coincides with the Ganges River peak in 2022. In contrast, sandy areas in the inlet of Ganges River are more pronounced in 2023 due to the changed conditions and local climate setting which is visible on the scenes of \hl{the classified }Landsat \hl{images}. The increase of flooded inundated areas in the wetlands of the Ganges River, the development of mangrove forests and the fragmentations of the forest land cover types were detected on the images for scenes processed \hl{using the} DL modules of \hl{the} GRASS GIS. The hydrological contribution from Meghna tributary to Ganges River stream follows a similar hydrological pattern to the one of Brahmaputra. Since Meghna has high water levels that extend into September due of back water effect, the contribution of Meghna to Ganges River flow is more pronounced on the images for November. Hence, as shown on the images, the categories of land cover types were visualised and evaluated using image classification methods. The flooding periods impacted the hydrology of the Ganges River by the monsoon effects in southern Bangladesh.

%----------- subsection ----------->
\subsection{K-means classification outcomes}

Flooded areas in Ganges River Delta detected and mapped using MaxLike methods are presented as a series of the resulting maps in Figure \ref{fig09} and Figure \ref{fig10} and deep learning (DL) is presented in Figure \ref{fig11} and Figure \ref{fig12}. 

\begin{figure}[H]
\hspace{50pt}\makebox[0.90\linewidth][r]{%
	\begin{subfigure}[b]{.40\textwidth}
		\centering
		\includegraphics[width=0.98\textwidth]{Fig_09a.jpg}
		\caption{2021}
	\end{subfigure}%
	\begin{subfigure}[b]{.40\textwidth}
		\centering
		\includegraphics[width=0.98\textwidth]{Fig_09b.jpg}
		\caption{2022}
	\end{subfigure}%
	\begin{subfigure}[b]{.40\textwidth}
		\centering
		\includegraphics[width=0.98\textwidth]{Fig_09c.jpg}
		\caption{2023}
	\end{subfigure}%
}
\caption{Flooded landscapes in Ganges River Delta, Bangladesh mapped using k-means clustering and MaxLike classification of GRASS GIS applied to the Landsat 8-9 OLI/TIRS images: \textbf{(a)}: March 2021; \textbf{(b)}: March 2022; \textbf{(c)}: March 2023. Ten land cover types correspond to the following categogies: 1); Water; 2) Wetlands and mudflats; 3) Mangrove forests; 4) Sandy areas; 5) Forests; 6) Croplands; 7) Grasslands; 8) Urban settlements; 9) Orchards; 10) Aquaculture.}\label{fig09}
\end{figure}

\begin{figure}[H]
\hspace{50pt}\makebox[0.90\linewidth][r]{%
	\begin{subfigure}[b]{.40\textwidth}
		\centering
			\includegraphics[width=0.98\textwidth]{Fig_10a.jpg}
		\caption{2021}
	\end{subfigure}%
	\begin{subfigure}[b]{.40\textwidth}
		\centering
		\includegraphics[width=0.98\textwidth]{Fig_10b.jpg}
		\caption{2022}
	\end{subfigure}%
	\begin{subfigure}[b]{.40\textwidth}
		\centering
		\includegraphics[width=0.98\textwidth]{Fig_10c.jpg}
		\caption{2023}
	\end{subfigure}%
}
\caption{Post-flood landscapes in Ganges River Delta, Bangladesh mapped using k-means clustering and MaxLike classification of GRASS GIS applied to the Landsat 8-9 OLI/TIRS images: \textbf{(a)}: November 2021; \textbf{(b)}: November 2022; \textbf{(c)}: November 2023. Ten land cover types correspond to the following categogies: 1); Water; 2) Wetlands and mudflats; 3) Mangrove forests; 4) Sandy areas; 5) Forests; 6) Croplands; 7) Grasslands; 8) Urban settlements; 9) Orchards; 10) Aquaculture.}\label{fig10}
\end{figure}

The classified images show the pre and post-flood inundation maps for the region of the Ganges River Delta during March and November on 2021, 2022 and 2023, based on the Landsat images. The findings revealed a changing pattern of expanded floods in the post-flood period in Ganges River Delta, and inequality in distribution of inundated areas covered by water by consequent years (2021, 2022 and 2023). 

The flood periods in 2022 in southern Bangladesh caused significant  deluge and inundation which affected cropland areas as well as  settlements in rural and urban districts. Besides, selected areas experienced extended periods of inundations, while coastal areas were continuously inundated. Finally, land cover types related to orchards were affected\hl{ by floods. }Thus, the results show the distribution of continuing waterbodies in March which increases to November covering an area of southern Bangladesh. In March 2021, a total flood-inundated area was lower, with most inundation was occurring in cropland (Class 6), followed by Urban settlement (Class 8) and orchard areas (Class 9), and other areas such as wetlands, mangroves and aquaculture (Classes 2, 3 and 10). 

The maps are generated using \hl{the} classified images for March showing the flooding event (in Figure \ref{fig09}) and Figure \ref{fig11}, as well as for November (Figure \ref{fig10} and Figure \ref{fig12}) showing the post-flood landscapes, respectively. Comparing these two periods, all crop-related land cover types present areas notably affected by floods, which also occurred in the regions of residential property. Correspondingly, larger areas were inundated with \hl{the} affected public infrastructure during the catastrophic November months when heavy monsoon rains continue to pummel the affected districts in \hl{a} close proximity to \hl{the} Ganges River. 

%----------- subsection ----------->
\subsection{Deep Learning outcomes}

The supervised learning models of GRASS GIS based on the MLPClassifier demonstrated accurate results of flood mapping and refined visualisation functionality\hl{. The outcomes of DL show} that for the March to November time frame, some patterns were common for both months, while autumn months were newly inundated, Figure \ref{fig11} and Figure \ref{fig12}. 

\begin{figure}[H]
\hspace{50pt}\makebox[0.90\linewidth][r]{%
	\begin{subfigure}[b]{.40\textwidth}
		\centering
			\includegraphics[width=0.98\textwidth]{Fig_11a.jpg}
		\caption{2021}
	\end{subfigure}%
	\begin{subfigure}[b]{.40\textwidth}
		\centering
		\includegraphics[width=0.98\textwidth]{Fig_11b.jpg}
		\caption{2022}
	\end{subfigure}%
	\begin{subfigure}[b]{.40\textwidth}
		\centering
		\includegraphics[width=0.98\textwidth]{Fig_11c.jpg}
		\caption{2023}
	\end{subfigure}%
}
\caption{Flooded landscapes in Ganges River Delta, Bangladesh mapped using Deep Learning (DL) classification MLPC of GRASS GIS applied to the Landsat 8-9 OLI/TIRS images: \textbf{(a)}: March 2021; \textbf{(b)}: March 2022; \textbf{(c)}: March 2023. Ten land cover types correspond to the following categogies: 1); Water; 2) Wetlands and mudflats; 3) Mangrove forests; 4) Sandy areas; 5) Forests; 6) Croplands; 7) Grasslands; 8) Urban settlements; 9) Orchards; 10) Aquaculture.}\label{fig11}
\end{figure}

In contrast, spring period demonstrated that selected northern regions recovered from the areas which experienced continuous inundation, while southern areas located in the proximity of Ganges River Delta were progressively inundated. While the \hl{DL} method was developed to serve a specific goal of image processing, its purpose is applicable as solutions to flood monitoring with \hl{the} outcomes of \hl{the} detected inundated areas\hl{, }as presented above. Hence, the research outcomes have provided \hl{valuable} insights into \hl{the monitoring of }changing landscapes \hl{comparing flooded and} post-flooded period of \hl{the} Ganges River Delta. 

\begin{figure}[H]
\hspace{50pt}\makebox[0.90\linewidth][r]{%
	\begin{subfigure}[b]{.40\textwidth}
		\centering
			\includegraphics[width=0.98\textwidth]{Fig_12a.jpg}
		\caption{2021}
	\end{subfigure}%
	\begin{subfigure}[b]{.40\textwidth}
		\centering
		\includegraphics[width=0.98\textwidth]{Fig_12b.jpg}
		\caption{2022}
	\end{subfigure}%
	\begin{subfigure}[b]{.40\textwidth}
		\centering
		\includegraphics[width=0.98\textwidth]{Fig_12c.jpg}
		\caption{2023}
	\end{subfigure}%
}
\caption{Post-flood landscapes in Ganges River Delta, Bangladesh mapped using Deep Learning (DL) classification MLPC of GRASS GIS applied to the Landsat 8-9 OLI/TIRS images: \textbf{(a)}: November 2021; \textbf{(b)}: November 2022; \textbf{(c)}: November 2023. Ten land cover types correspond to the following categogies: 1); Water; 2) Wetlands and mudflats; 3) Mangrove forests; 4) Sandy areas; 5) Forests; 6) Croplands; 7) Grasslands; 8) Urban settlements; 9) Orchards; 10) Aquaculture.}\label{fig12}
\end{figure}

The time-series analysis of the Ganges River Delta was used to plots and graphically display\hl{ }changes in the flood and post-flood periods for southern Bangladesh\hl{. The analysis of data }for three test periods (2021, 2022 and 2023) illustrated some recovered from the flood waters as time progressed which contrasted with the inundated areas. The data for 2014 were used as training polygon for \hl{DL} modelling which requires the classification seed. The derived maps were produced from the cloud-free Landsat images processed by the \hl{DL}  methods between the flooded period in March, and post-flooded period in November where some areas recovered from the flood inundation. The maps generated using GRASS GIS present show high potential for their adoption in environmental monitoring. The methods of deep learning supported accurate image classification through the advanced computer vision algorithms of pattern recognition. 

\begin{table}[H]
\caption{Flood period: estimated classes of land cover types for March period for years 2021, 2022 and 2023 in Ganges River Delta}
  \centering
  \begin{tabularx}{\textwidth}{p{0.9cm}p{0.9cm}p{0.9cm}p{0.9cm}p{0.9cm}p{0.9cm}p{0.9cm}p{0.9cm}p{0.9cm}p{0.9cm}p{0.9cm}}
  \toprule
  \multirow{2}{*}{Year} & \multicolumn{10}{c}{Classes of land cover types in March pixels: Ganges-Brahmaputra River Delta} \\
     & 1 & 2 & 3 & 4 & 5 & 6 & 7 & 8 & 9 & 10 \\
    \midrule
    \textbf{2021} & 717 & 252 & 434 & 758 & 1005 & 1035 & 758 & 744 & 927 & 219 \\
    \textbf{2022} &  749 & 236 & 489 & 803 & 1030 & 1100 & 557 & 744 & 800 & 336 \\
     \textbf{2023} &  707 & 343 & 1055 & 381 & 1032 & 997 & 702 & 649 & 796 & 180 \\
\bottomrule
\end{tabularx}
  \label{tab02}
\end{table}

\begin{table}[H]
 \caption{Post-flood period: estimated classes of land cover types for November period for years 2021, 2022 and 2023 in Ganges River Delta}
  \centering
  \begin{tabularx}{\textwidth}{p{0.9cm}p{0.9cm}p{0.9cm}p{0.9cm}p{0.9cm}p{0.9cm}p{0.9cm}p{0.9cm}p{0.9cm}p{0.9cm}p{0.9cm}}
  \toprule
  \multirow{2}{*}{Year} & \multicolumn{10}{c}{Classes of land cover types in November pixels: Ganges-Brahmaputra River Delta} \\
     & 1 & 2 & 3 & 4 & 5 & 6 & 7 & 8 & 9 & 10 \\
    \midrule
    \textbf{2021} & 689 & 402 & 579 & 565 & 1228 & 1109 & 571 & 849 & 647 & 208 \\
    \textbf{2022} & 700 & 468 & 365 & 587 & 1066 & 1048 & 683 & 942 & 683 & 326 \\
     \textbf{2023} & 715 & 511 & 231 & 599 & 1351 & 1132 & 718 & 707 & 685 & 310 \\
\bottomrule
\end{tabularx}
  \label{tab03}
\end{table}

The time series flood data made using scripting models of GRASS GIS with open source available specifications catalyse further innovation in risk management and presents solutions for state-of-the-art mapping of floods using programming approaches and \hl{ML} modules. The following 10 land cover types were identified around the Ganges River Delta, Bangladesh: \emph{1); Water; 2) Wetlands and mudflats; 3) Mangrove forests; 4) Sandy areas; 5) Forests; 6) Croplands; 7) Grasslands; 8) Urban settlements; 9) Orchards; 10) Aquaculture.} Information on land cover types is derived from Rahman et al. \cite{Rahman2020}, Sousa and Small \cite{rs13234799} and Paszkowski et al. \cite{Paszkowski2021} and applied to the extent of the current study area. The presented maps of floods can be used as a broad base of environmental applications and flood hazard assessment. Their application can range from academia and research, industry and development, to governmental and social services for risk prevention. 

The maps derived using deep learning approach were produced from the cloud-free Landsat scenes between March and November periods and are shown in Figures \ref{fig11} and \ref{fig12}. The presented land cover maps are intended for potential flood damage assessment in southern Bangladesh corrections of the presented framework and Python-based algorithm of GRASS GIS using machine learning (ML) approach which uses MLPClassifier of Deep Learning (DL) presents an advanced tool for automated recognition of land cover types on the satellite images and monitoring floods in the prone to floods regions of the southern Bangladesh. The percentage of the correctly classified images was calculated and summarised in the matrices showing classes separability (Matrix \ref{eq1}, Matrix \ref{eq2} and Matrix \ref{eq3} for years 2021, 2022 and 2023, respectively) and evaluated in the accuracy assessment. 

Class separability matrices of land cover classes identified over Ganges River Delta which are presented in the \hl{Figures }\ref{fig13}, \ref{fig14} and \ref{fig15} \hl{which show the following} matrices below: Matrix \ref{eq1} shows class separability matrices for March and November 2021; Matrix \ref{eq2} shows class separability matrices for March and November 2022, and Matrix \ref{eq3} shows class separability matrices for March and November 2023\hl{, respectively}. The covariance matrices signify the values of the class separability which are composed of the cluster means based on the values of the spectral signatures of pixels. They demonstrate the results of the cluster's means \hl{obtained from the computed} the covariance matrices\hl{. The matrices} are generated \hl{using} the 'i.maxlik' function of GRASS GIS during the image classification process. The columns \hl{signify} 10 generated land cover classes, and the rows signify the number of pixels which are assigned to those classes correctly. Hence, the class separability matrices show the error matrix for \hl{the} pixel-based maximum-likelihood classification which indicate the correctness of the pixels assigned to the \hl{land cover }classes from 1 to 10. The result of \hl{image} clustering\hl{ }using values of spectral classes are related to \hl{the} land cover types which are identified on the ground of the Ganges River Delta and \hl{the} surroundings. 

The matrices consist of ten classes and computed pixels for each land cover class. The land categories are defined as follows: waterbodies, wetlands and mudflats indicating inundated areas, mangrove forests and associated swamps, sandy areas in the riverbanks, forests (e.g., tree cover, hill forests); croplands and homestead croplands, grasslands, urban and rural settlements, orchards and aquaculture. Class separability matrices evaluated the assignment of the pixels distributed to various land cover classes over the landscapes of Ganges River Delta. Hence the matrices serve as a function of the regularisation parameters used for distinguishability of pixels according to their spectral reflectances. In this way, class separability matrices perform \hl{the} iterative search for cells within the raster matrix of the image considering a range of values of pixels assigned to different land cover types.

\begin{figure}[H]
\begin{equation}\label{eq1}
\setlength\arraycolsep{3pt}
\footnotesize
\begin{vmatrix}
    & 1 & 2 & 3 & 4 & 5 & 6 & 7 & 8 & 9 & 10 \\
1 & 0 &  &  &  &  &  &  &  &  &  \\
2 & 1.3 & 0 &  &  &  &  &  &  &  &  \\
3 & 2.2 & 1.0 & 0 &  &  &  &  &  &  &  \\
4 & 3.0 & 1.5 & 0.7 & 0 &  &  &  &  &  &  \\
5 & 3.2 & 1.2 & 1.2 & 1.2 & 0 &  &  &  &  &  \\
6 & 3.4 & 1.9 & 1.3 & 0.7 & 1.1 & 0 &  &  &  &  \\  
7 & 3.6 & 2.4 & 1.7 & 1.2 & 1.6 & 0.7 & 0 &  &  &  \\  
8 & 3.7 & 1.7 & 1.8 & 1.7 & 0.7 & 1.4 & 1.7 & 0 &  &  \\
9 & 3.4 & 1.6 & 2.3 & 2.5 & 1.3 & 2.3 & 2.6 & 1.0 & 0 &  \\   
10 & 3.8 & 2.2 & 2.8 & 2.8 & 2.0 & 2.7 & 3.0 & 1.8 & 1.1 & 0 \\      
\end{vmatrix}
 %
 \hspace{10pt}
 \begin{vmatrix}
    & 1 & 2 & 3 & 4 & 5 & 6 & 7 & 8 & 9 & 10 \\
1 & 0 &  &  &  &  &  &  &  &  &  \\
2 & 1.9 & 0 &  &  &  &  &  &  &  &  \\
3 & 2.6 & 1.0 & 0 &  &  &  &  &  &  &  \\
4 & 2.9 & 2.2 & 1.4 & 0 &  &  &  &  &  &  \\
5 & 3.9 & 2.0 & 0.9 & 1.2 & 0 &  &  &  &  &  \\
6 & 4.6 & 2.4 & 1.4 & 1.5 & 0.6 & 0 &  &  &  &  \\  
7 & 4.2 & 2.8 & 1.9 & 1.3 & 1.3 & 1.0 & 0 &  &  &  \\  
8 & 4.1 & 2.9 & 1.8 & 0.8 & 1.3 & 1.3 & 0.7 & 0 &  &  \\
9 & 3.9 & 3.3 & 2.5 & 1.0 & 2.2 & 2.3 & 1.5 & 1.1 & 0 &  \\   
10 & 3.5 & 3.3 & 2.6 & 1.5 & 2.3 & 2.3 & 1.8 & 1.5 & 0.9 & 0 \\       
 \end{vmatrix} 
\end{equation}
\caption{\hl{Class separability matrices computed for classified Landsat 8-9 OLI/TIRS satellite images on March and November 2021}}
\label{fig13}
\end{figure}

\begin{figure}[H]
\begin{equation}\label{eq2}
\setlength\arraycolsep{3pt}
\footnotesize
\begin{vmatrix}
   & 1 & 2 & 3 & 4 & 5 & 6 & 7 & 8 & 9 & 10 \\
1 & 0 &  &  &  &  &  &  &  &  &  \\
2 & 1.5 & 0 &  &  &  &  &  &  &  &  \\
3 & 2.4 & 1.0 & 0 &  &  &  &  &  &  &  \\
4 & 3.3 & 1.7 & 0.7 & 0 &  &  &  &  &  &  \\
5 & 3.4 & 1.3 & 1.3 & 1.3 & 0 &  &  &  &  &  \\
6 & 3.8 & 2.2 & 1.3 & 0.7 & 1.2 & 0 &  &  &  &  \\  
7 & 3.6 & 2.4 & 1.6 & 1.1 & 1.5 & 0.6 & 0 &  &  &  \\  
8 & 3.9 & 1.8 & 1.9 & 1.8 & 0.7 & 1.4 & 1.6 & 0 &  &  \\
9 & 3.9 & 1.7 & 2.5 & 2.7 & 1.4 & 2.4 & 2.5 & 1.1 & 0 &  \\   
10 & 3.9 & 2.1 & 2.8 & 2.9 & 2.0 & 2.7 & 2.8 & 1.7 & 1.0 & 0 \\      
\end{vmatrix}
 %
 \hspace{10pt}
 \begin{vmatrix}
    & 1 & 2 & 3 & 4 & 5 & 6 & 7 & 8 & 9 & 10 \\
1 & 0 &  &  &  &  &  &  &  &  &  \\
2 & 1.8 & 0 &  &  &  &  &  &  &  &  \\
3 & 2.4 & 1.8 & 0 &  &  &  &  &  &  &  \\
4 & 2.7 & 0.9 & 1.3 & 0 &  &  &  &  &  &  \\
5 & 3.8 & 1.8 & 1.4 & 0.8 & 0 &  &  &  &  &  \\
6 & 4.5 & 2.2 & 1.8 & 1.3 & 0.6 & 0 &  &  &  &  \\  
7 & 4.1 & 2.5 & 1.5 & 1.7 & 1.3 & 1.0 & 0 &  &  &  \\  
8 & 4.1 & 2.5 & 1.0 & 1.5 & 1.1 & 1.2 & 0.7 & 0 &  &  \\
9 & 3.9 & 2.9 & 0.8 & 2.3 & 2.2 & 2.3 & 1.4 & 1.2 & 0 &  \\   
10 & 3.0 & 2.8 & 1.2 & 2.3 & 2.3 & 2.4 & 1.8 & 1.7 & 0.8 & 0 \\       
 \end{vmatrix} 
\end{equation}
\caption{\hl{Class separability matrices computed for classified Landsat 8-9 OLI/TIRS satellite images on March and November 2022}}
\label{fig14}
\end{figure}

\begin{figure}[H]
\begin{equation}\label{eq3}
\setlength\arraycolsep{3pt}
\footnotesize
\begin{vmatrix}
    & 1 & 2 & 3 & 4 & 5 & 6 & 7 & 8 & 9 & 10 \\
1 & 0 &  &  &  &  &  &  &  &  &  \\
2 & 1.4 & 0 &  &  &  &  &  &  &  &  \\
3 & 3.0 & 1.0 & 0 &  &  &  &  &  &  &  \\
4 & 3.0 & 1.3 & 1.4 & 0 &  &  &  &  &  &  \\
5 & 3.8 & 1.7 & 0.7 & 1.6 & 0 &  &  &  &  &  \\
6 & 4.1 & 1.7 & 1.2 & 0.9 & 1.1 & 0 &  &  &  &  \\  
7 & 3.9 & 2.1 & 1.3 & 2.1 & 0.7 & 1.5 & 0 &  &  &  \\  
8 & 4.4 & 2.2 & 1.7 & 1.2 & 1.4 & 0.7 & 1.4 & 0 &  &  \\
9 & 4.2 & 2.3 & 2.3 & 1.1 & 2.3 & 1.4 & 2.4 & 1.1 & 0 &  \\   
10 & 3.9 & 2.6 & 2.6 & 1.8 & 2.5 & 2.0 & 2.6 & 1.7 & 1.1 & 0 \\      
\end{vmatrix}
 %
 \hspace{10pt}
 \begin{vmatrix}
    & 1 & 2 & 3 & 4 & 5 & 6 & 7 & 8 & 9 & 10 \\
1 & 0 &  &  &  &  &  &  &  &  &  \\
2 & 1.8 & 0 &  &  &  &  &  &  &  &  \\
3 & 2.1 & 1.4 & 0 &  &  &  &  &  &  &  \\
4 & 3.1 & 1.1 & 1.2 & 0 &  &  &  &  &  &  \\
5 & 3.9 & 1.9 & 1.3 & 0.8 & 0 &  &  &  &  &  \\
6 & 4.2 & 2.3 & 1.7 & 1.2 & 0.6 & 0 &  &  &  &  \\  
7 & 3.5 & 2.5 & 0.9 & 1.8 & 1.4 & 1.5 & 0 &  &  &  \\  
8 & 4.1 & 2.6 & 1.7 & 1.8 & 1.3 & 0.9 & 1.2 & 0 &  &  \\
9 & 4.0 & 2.9 & 1.4 & 2.2 & 1.7 & 1.6 & 0.6 & 0.9 & 0 &  \\   
10 & 3.1 & 2.7 & 1.6 & 2.2 & 1.9 & 1.9 & 1.0 & 1.6 & 0.9 & 0 \\       
 \end{vmatrix} 
\end{equation}
\caption{\hl{Class separability matrices computed for classified Landsat 8-9 OLI/TIRS satellite images on March and November 2023}}
\label{fig15}
\end{figure}

Through the applied automatic DL-based image processing and mapping workflow, flood inundation mapping of southern Bangladesh \hl{has} the potential to provide information on inundated areas\hl{. Time series of Landsat images are essential }for flood management in real time regime. Using the techniques of deep learning applied for flood mapping, a significant difference in values of spectral reflectance on land cover classes for inundated areas (containing high percentage of water and moisture) and non-water (dry) areas enables a distinct separation between these land cover categories. Automatically derived ranges for inundated areas in Landsat images show a clear discrimination of \hl{the} affected areas from other classes, as presented in \hl{the }figures \hl{supporting this research}. 

It is well known that Bangladesh is a country particularly vulnerable to \hl{the} hydrological hazards and flood disasters. Therefore, testing and implementing novel cartographic methods for RS-based mapping of flooded areas is essential for operative monitoring of areas at risks, such as coastal areas of Ganges River Delta. The cumulative effects of flat topographic relief and monsoon climate resulted in high vulnerability of Bangladesh to the floods with over 80\% of the population is exposed to flood risks. Diverse measures are undertaken by the authorities and government of this coastal country to protect population from the devastating flood disasters. \hl{The} engineering constructions against floods include, for instance, embankments, flood shelters, hydrologic sluice gates and river regulators, dredged drainage channels and soon. 

\hl{Earth} observation data present\hl{ indispensable} resources for flood monitoring from space \hl{and can be used for flood monitoring, as demonstrated in this study}. \hl{Technical} advantage of the demonstrated novel methods of flood mapping consists in automated approach of deep learning applied for inundated areas in Bangladesh. \hl{More specifically}, this paper contributed to the development of cartographic methods for visualisation of flooded areas using satellite images and deep learning (DL) techniques available by GRASS GIS. The presented and commented programming codes can be reused \hl{for} flood mapping in similar studies. 

%----------- section ----------->
\section{Discussion}

Floods are a major threat to human communities all over the world, with serious socio-economic and environmental consequences, given that they are the leading cause of natural disaster losses. Related research shows that one in four people worldwide is considered to be at significant risk from flooding. Bangladesh, which occupies low-lying flood and tidal plains, has one of the world's largest and most disaster-prone deltas. \hl{In particular}, the absolute number of people at risk of flooding \hl{in Bangladesh} is 94.42 million, or 57.5\% of the country's population (the second highest in the world). \hl{Consequently, }floods in Ganges River Delta present one of the most catastrophic problems at the national level in \hl{the country} \cite{Szabo}. With this regards, preventive mapping of flooded areas presents the societal importance for the social system of Bangladesh and \hl{contributes to }maintaining the sustainability of both natural landscapes and social infrastructure. 

Social awareness of the flood hazards can help mitigating the risk associated with floods. \hl{Remote} sensing (RS) data processed using advanced methods and visualised on maps \hl{are} essential \hl{for} addressing \hl{the} environmental challenges\hl{. Processing RS data is useful for} mapping the areas at risks, cartographic display of the inundated coastal and prognosis of flood in deltaic areas based on the time series analysis. To this end, the advanced methods of deep learning (DL) used for operative mapping of flooded areas support publicly accessible, modifiable, and distributable open-source cartographic products. Such products can be used for preventive monitoring of the areas at risk, prognosis of flood events and hydrologic modelling. Hence, operative monitoring and forecasting of flood hazards \hl{using RS data }can effectively reduce the catastrophic consequences of floods in Ganges River Delta. 

Advanced methods of \hl{DL} for accurate processing of satellite images are possible using open-source GRASS GIS which presents an effective method to retrieve geo-information. This study presented \hl{the use of }such DL-based image analysis using\hl{ }GRASS GIS \hl{for} mapping flooded areas in the coastal region of the Ganges River Delta\hl{. }By comparison of \hl{the} multi-temporal RS data, the areas affected by floods were detected and compared for the flood- and post-flood periods. \hl{The} presented workflow demonstrated that DL techniques of GRASS GIS \hl{have} significant advantages for image classification \hl{because the DL increases} the accuracy and automation of mapping through the algorithms of AI employed in the cartographic workflow. Among others,\hl{ }the automated detection of similar spectral characteristics \hl{was applied to} different \hl{land categories} and assigning them to distinct classes which enables to overcome the challenges of the misclassification. Second, the heterogeneous landscapes with highly dense mosaic of landscape \hl{patterns}, typical for coastal Bangladesh, are identified by the ML\hl{ }with regard to\hl{ }resolution of the Landsat images and defined classes. Spectral properties of the multispectral satellite images\hl{ }were evaluated \hl{for} land cover changes\hl{ }in the inundated areas of the Ganges River Delta\hl{.} 

Flooded areas in the Ganges River Delta were determined over the target dates using time series analysis of the satellite images Landsat, and compared for flood- and post-flood periods. The results indicated the increase in wetlands and inundated areas during the later period detected by the DL methods\hl{ }and cartographic scripts. The obtained results presented a series of maps on the Ganges River Delta that showed that flood periods reached a peak in November and affected coastal landscapes of the Ganges River, and demonstrated the increase in flooded areas. The presented time series of maps showing floods in March and post-flooded landscapes in November highlighted the benefits achieved by \hl{RS} data visualization for analysis \hl{and comparison of} flood extent\hl{.} 

Floods have varied effect on the behaviour of landscapes which differed over years, as indirect indicator of climate effects and monsoon processes. In fact, the combination of regional hydrologic behaviour of Ganges River, the structure of sediment, the proportions of clayey sediments in the soil results in the degree of water stagnation with significantly increased the extent of inundated areas. The behaviour of coastal areas was evaluated within three years (for 2021, 2022 and 2023) to demonstrate that the effects of floods on landscape which considerably altered the landscapes as it is reflected on the maps prepared using clustering (Figures \ref{fig09}-\ref{fig10}) and DL techniques (Figures \ref{fig11}-\ref{fig12}), as evaluated during 3 years. Similar to existing studies discussed in earlier sections, the flooded period of the Ganges River resulted in the highest level of inundation of the coastal regions of southern Bangladesh in November period compared to March. 

\section{Conclusions}
This study demonstrated the capacity of sequential analysis of the RS data to identify changes in the coastal and wetland landscapes of Ganges River Delta over the last three years\hl{ }for comparison of flood event with post-flood landscapes. The results illustrated the effectiveness of the satellite images in identification of the flood extent and detecting the affected areas. Moreover, the results shown the technological advantages of the DL for satellite image processing with a case of GRASS GIS. \hl{The} demonstrated and commented codes can be used in future similar studies on vulnerability\hl{analysis} and flood risk assessment. In relation to \hl{the} previous studies and\hl{ }working hypotheses, the results of \hl{this} study can be integrated with related \hl{geoinformation}, such as hydrologic models, satellite altimetry, aerial imagery\hl{,} topographic maps, \hl{and }hydrodynamic flood models. These data can be then be used for flood forecasting, simulating flood models in southern Bangladesh, providing information on spatial extent of the flooded area for estimating possible impact of flood. Integrated use of such information is useful  for undertaking adaptive measures and mitigating floods in Ganges River Delta.

Importantly, this paper presented a case of flood mapping in Ganges River Delta using the integration of \hl{the} open source data and methods: the dataset was obtained from the publicly available repository of the USGS and processed using freely available GRASS GIS software. Such approach is essential for\hl{ }policy makers from the developing countries such as Bangladesh, since free data and tools can motivate further studies to continue and expand spatial analysis using the proposed techniques. In future directions of the environmental monitoring, this research can be further extended for larger spatio-temporal extent of the Ganges River Delta\hl{. For example, it can be done }by producing a series of maps covering extended periods using complementary RS data. Moreover, upscaling the data for a basin or sub-basin level can support more general estimation of flood impacts for hydrological risk assessment. Finally, spatial analysis based on ML-based image processing can be included as a cartographic support into the community-based water resource management of southern Bangladesh to monitor areas prone to flood disasters.

\hl{To conclude, }this research has demonstrated an example of using the advanced approach of ML in cartography and RS data processing. Thus, it has led to a greater understanding of the importance of using advanced methods of machine learning for environmental monitoring. The application of the ANN as a branch of machine learning methods, presented here as a case of MLP classifier algorithm demonstrated that such technique can be successfully applied in similar studies which require RS data for spatio-temporal analysis. Using ML methods of GRASS GIS remote sensing software demonstrated a fine perspectives for image processing and environmental data analysis. We evaluated the proposed ML framework of GRASS GIS on the Landsat scenes, however, other satellite images can also be applied in future similar studies.
%----------- section ----------->

\institutionalreview{Not applicable.}

\informedconsent{Not applicable.}

\dataavailability{The raw data supporting the conclusions of this article will be made available by the authors on request. The binary results of the DL mapping are available in the GitHub repository, including GRASS GIS scripts and the results of the DL image analysis: \href{https://github.com/paulinelemenkova/Floods_Ganges_GRASS_GIS_DL_Image_Analysis}{https://github.com/paulinelemenkova/Floods\_Ganges\_GRASS\_GIS\_DL\_Image\_Analysis} (accessed on 5 April 2024).}

\acknowledgments{The author thanks the reviewers for reading and review of this manuscript.}

\conflictsofinterest{The author declares that they have no known competing financial interests neither personal relationships that could have influenced the work reported in this article.} 

\abbreviations{Abbreviations}{
The following abbreviations are used in this manuscript:\\

\noindent 
\begin{tabular}{@{}ll}
\hl{AI} & \hl{Artificial Intelligence} \\
ANN & Artificial Neural Networks \\
DL & Deep Learning \\
DN & Digital Number \\
EO & Earth Observation \\
GEBCO & General Bathymetric Chart of the Oceans \\
GMT & Generic Mapping Tools \\
GRASS & Geographic Resources Analysis Support System \\
GIS & Geographic Information System \\
Landsat OLI/TIRS & Landsat Operational Land Imager and Thermal Infrared Sensor \\
ML & Machine Learning \\
MLP & MultiLayer Perceptron \\
NIR & Near Infrared \\ 
RF & Random Forest \\
RS & Remote Sensing \\
SVM & Support Vector Machines \\
SWIR & Shortwave Infrared \\
USGS & United States Geological Survey \\
UTM & Universal Transverse Mercator \\
WGS84 & World Geodetic System 84\\
WRS & Worldwide Reference System  
\end{tabular}
}

\appendixtitles{no} % Leave argument "no" if all appendix headings stay EMPTY (then no dot is printed after "Appendix A"). If the appendix sections contain a heading then change the argument to "yes".
\appendixstart
\appendix
\section[\appendixname~\thesection]{Metadata for satellite images }

\begin{table}[H]
\footnotesize
\caption{Metadata for the Landsat 8-9 OLI/TIRS images obtained from the USGS} \label{tabApp}
\begin{adjustwidth}{-\extralength}{0cm}
    \begin{tabularx}{\fulllength}{llll}
    \toprule
       	\textbf{Data Set Attribute} & \textbf{Attribute Value} & \textbf{Attribute Value} & \textbf{Attribute Value} \\ \hline
        \textbf{Landsat Scene Identifier} & LC81370442023066LGN00 & LC81370442022079LGN00 & LC81370442021076LGN00 \\ \hline
        \textbf{Date Acquired} & \textbf{2023/03/07} & \textbf{2022/03/20} & \textbf{2021/03/17} \\ \hline
        \textbf{Roll Angle} & -1 & 0 & 0 \\ \hline
        \textbf{Start Time} & 2023-03-07 04:24:33 & 2022-03-20 04:24:26.058914 & 2021-03-17 04:24:23.120295 \\ \hline
        \textbf{Stop Time} & 2023-03-07 04:25:05 & 2022-03-20 04:24:57.828914 & 2021-03-17 04:24:54.890294 \\ \hline
        \textbf{Land Cloud Cover} & 3.12 & 0.01 & 0.03 \\ \hline
        \textbf{Scene Cloud Cover L1} & 2.97 & 0.01 & 0.03 \\ \hline
        \textbf{Ground Control Points Model} & 732 & 771 & 776 \\ \hline
        \textbf{Ground Control Points Version} & 5 & 5 & 5 \\ \hline
        \textbf{Geometric RMSE Model} & 5683 & 5615 & 5694 \\ \hline
        \textbf{Geometric RMSE Model X} & 3887 & 3857 & 3585 \\ \hline
        \textbf{Geometric RMSE Model Y} & 4146 & 4081 & 4424 \\ \hline
        \textbf{Processing Software Version} & LPGS\_16.2.0 & LPGS\_15.6.0 & LPGS\_15.4.0 \\ \hline
        \textbf{Sun Elevation L0RA} & 51.70012312 & 56.10505202 & 55.19368850 \\ \hline
        \textbf{Sun Azimuth L0RA} & 134.86314148 & 129.90704446 & 131.01649112 \\ \hline
        \textbf{TIRS SSM Model} & FINAL & FINAL & FINAL \\ \hline
        \textbf{Data Type L2} & OLI\_TIRS\_L2SP & OLI\_TIRS\_L2SP & OLI\_TIRS\_L2SP \\ \hline
        \textbf{Satellite} & 8 & 8 & 8 \\ \hline
        \textbf{Scene Center Lat DMS} & "23°06'46.58""N" & "23°06'45.29""N" & "23°06'46.01""N" \\ \hline
        \textbf{Scene Center Long DMS} & "90°23'29.83""E" & "90°23'14.57""E" & "90°24'53.42""E" \\ \hline
        \textbf{Corner Upper Left Lat DMS} & "24°09'00.29""N" & "24°09'00""N" & "24°09'02.38""N" \\ \hline
        \textbf{Corner Upper Left Long DMS} & "89°13'54.73""E" & "89°13'44.11""E" & "89°15'19.58""E" \\ \hline
        \textbf{Corner Upper Right Lat DMS} & "24°11'21.62""N" & "24°11'21.52""N" & "24°11'22.45""N" \\ \hline
        \textbf{Corner Upper Right Long DMS} & "91°30'23.18""E" & "91°30'12.56""E" & "91°31'48.22""E" \\ \hline
        \textbf{Corner Lower Left Lat DMS} & "22°01'28.45""N" & "22°01'28.20""N" & "22°01'30.32""N" \\ \hline
        \textbf{Corner Lower Left Long DMS} & "89°17'26.70""E" & "89°17'16.22""E" & "89°18'50.22""E" \\ \hline
        \textbf{Corner Lower Right Lat DMS} & "22°03'36.04""N" & "22°03'35.93""N" & "22°03'36.76""N" \\ \hline
        \textbf{Corner Lower Right Long DMS} & "91°31'47.39""E" & "91°31'36.95""E" & "91°33'11.12""E" \\ \hline
        \midrule
        \textbf{Data Set Attribute} & \textbf{Attribute Value} & \textbf{Attribute Value} & \textbf{Attribute Value} \\ \hline
        \textbf{Landsat Scene Identifier} & LC91370442023330LGN00 & LC91370442022327LGN01 & LC81370442021332LGN00 \\ \hline
        \textbf{Date Acquired} & \textbf{2023/11/26} & \textbf{2022/11/23} & \textbf{2021/11/28} \\ \hline
        \textbf{Roll Angle} & 0 & 0 & 0 \\ \hline
        \textbf{Start Time} & 2023-11-26 04:24:51 & 2022-11-23 04:25:01 & 2021-11-28 04:24:54.925489 \\ \hline
        \textbf{Stop Time} & 2023-11-26 04:25:23 & 2022-11-23 04:25:32 & 2021-11-28 04:25:26.695489 \\ \hline
        \textbf{Land Cloud Cover} & 0.02 & 0.20 & 0.40 \\ \hline
        \textbf{Scene Cloud Cover L1} & 0.02 & 0.19 & 0.38 \\ \hline
        \textbf{Ground Control Points Model} & 750 & 782 & 768 \\ \hline
        \textbf{Ground Control Points Version} & 5 & 5 & 5 \\ \hline
        \textbf{Geometric RMSE Model} & 6651 & 6490 & 6697 \\ \hline
        \textbf{Geometric RMSE Model X} & 4185 & 4191 & 4329 \\ \hline
        \textbf{Geometric RMSE Model Y} & 5170 & 4955 & 5110 \\ \hline
        \textbf{Processing Software Version} & LPGS\_16.3.1 & LPGS\_16.2.0 & LPGS\_15.5.0 \\ \hline
        \textbf{Sun Elevation L0RA} & 41.83496249 & 42.43660966 & 41.36291892 \\ \hline
        \textbf{Sun Azimuth L0RA} & 154.44654325 & 154.42300606 & 154.52745495 \\ \hline
        \textbf{TIRS SSM Model} & N/A & N/A & FINAL \\ \hline
        \textbf{Data Type L2} & OLI\_TIRS\_L2SP & OLI\_TIRS\_L2SP & OLI\_TIRS\_L2SP \\ \hline
        \textbf{Satellite} & 9 & 9 & 8 \\ \hline
        \textbf{Scene Center Lat DMS} & "23°06'45.65""N" & "23°06'46.48""N" & "23°06'45.22""N" \\ \hline
        \textbf{Scene Center Long DMS} & "90°22'20.75""E" & "90°23'11.94""E" & "90°24'08.21""E" \\ \hline
        \textbf{Corner Upper Left Lat DMS} & "24°08'48.98""N" & "24°08'50.28""N" & "24°09'01.33""N" \\ \hline
        \textbf{Corner Upper Left Long DMS} & "89°12'51.37""E" & "89°13'44.40""E" & "89°14'37.14""E" \\ \hline
        \textbf{Corner Upper Right Lat DMS} & "24°11'11.26""N" & "24°11'11.80""N" & "24°11'22.06""N" \\ \hline
        \textbf{Corner Upper Right Long DMS} & "91°29'19.50""E" & "91°30'12.67""E" & "91°31'05.70""E" \\ \hline
        \textbf{Corner Lower Left Lat DMS} & "22°01'27.01""N" & "22°01'28.20""N" & "22°01'19.67""N" \\ \hline
        \textbf{Corner Lower Left Long DMS} & "89°16'24.02""E" & "89°17'16.22""E" & "89°18'08.71""E" \\ \hline
        \textbf{Corner Lower Right Lat DMS} & "22°03'35.46""N" & "22°03'35.93""N" & "22°03'26.64""N" \\ \hline
        \textbf{Corner Lower Right Long DMS} & "91°30'44.60""E" & "91°31'36.95""E" & "91°32'29.36""E" \\ \hline
        \bottomrule
    \end{tabularx}
    \end{adjustwidth}
\end{table}

%%%%%%%%%%%%%%%%%%%%%%%%%%%%%%%%%%%%%%%%%%
\begin{adjustwidth}{-\extralength}{0cm}
%\printendnotes[custom] % Un-comment to print a list of endnotes

\reftitle{References}
%\nocite{*}

%=====================================
% References, variant A: external bibliography
%=====================================
\bibliography{Ganges.bib}

%%%%%%%%%%%%%%%%%%%%%%%%%%%%%%%%%%%%%%%%%%
\end{adjustwidth}
\end{document}
